\usepackage[margin=0.65in]{geometry}
\usepackage[T1]{fontenc}
\usepackage{longtable}
\usepackage{booktabs}
\usepackage{graphicx}
\usepackage{xcolor}
\usepackage{listings}
\usepackage{array}
\usepackage{amsmath}
\usepackage{microtype}
\usepackage{needspace}
\usepackage[hypcap=false]{caption}
\usepackage{hyperref}
\usepackage{makeidx}
\makeindex
\setlength{\tabcolsep}{3.5pt}
\renewcommand{\arraystretch}{1.08}
\newcolumntype{L}[1]{>{\raggedright\arraybackslash}p{#1}}
\newcolumntype{R}[1]{>{\raggedleft\arraybackslash}p{#1}}
\newcolumntype{C}[1]{>{\centering\arraybackslash}p{#1}}
\lstdefinestyle{ptocode}{basicstyle=\ttfamily\tiny,breaklines=true,breakatwhitespace=false,columns=fullflexible,keepspaces=true,showstringspaces=false,frame=none,xleftmargin=0.75em,xrightmargin=0.25em,aboveskip=4pt,belowskip=8pt}
\lstdefinestyle{ptoscalarcode}{basicstyle=\ttfamily\small,breaklines=true,breakatwhitespace=false,columns=fullflexible,keepspaces=true,showstringspaces=false,frame=none,xleftmargin=0.75em,xrightmargin=0.25em,aboveskip=3pt,belowskip=5pt}
\setlength{\LTpre}{3pt}
\setlength{\LTpost}{6pt}
\newcommand{\TileInstructionEncoding}[5]{%
\Needspace{0.42\textheight}
{\small
\noindent\textbf{Encoding overview.} A PTO tile instruction is a variable-length descriptor parcel. Word 0 is always \texttt{B.OP}; \texttt{B.OP.desc\_count} declares how many following \texttt{B.*} descriptor words belong to the same CPU-visible instruction. The first encoding diagram shows the descriptor parcel order for this instruction. The second encoding diagram shows the \texttt{B.OP} bit allocation that selects the catalog opcode, ROM entry, descriptor count, variant, and profile-local selector.\par\smallskip
\noindent\includegraphics[width=\linewidth]{#1}\par\smallskip
\noindent\includegraphics[width=\linewidth]{#2}\par\smallskip
\begin{tabular}{@{}L{0.21\linewidth}L{0.19\linewidth}L{0.52\linewidth}@{}}
\toprule
Encoding field & Value & Meaning \\
\midrule
\texttt{Tile catalog opcode} & \texttt{#3} & PTO tile instruction catalog value. \\
\texttt{Descriptor bundle} & \texttt{#5} & Ordered CPU-visible descriptor bundle for this instruction profile. \\
\texttt{B.OP.OpcodeSel} & profile-defined & Profile-local selector carried in \texttt{B.OP}; combined with descriptor and extension context to identify this catalog entry. \\
\texttt{B.OP.desc\_count} & profile-defined & Number of following \texttt{B.*} descriptor instructions required by this instruction family. \\
\texttt{B.OP.rom\_entry\_id} & \texttt{#4} & Microcode ROM body selected after descriptor decode. \\
\bottomrule
\end{tabular}
}\par\smallskip
\noindent\textbf{Microcode encoding.} Each site in the table below is one fixed 64-bit vector micro instruction. The low 32 bits are \texttt{WM}; the high 32 bits are \texttt{WO}. The \texttt{WO} and \texttt{WM} columns list the two encoded halves of the same micro instruction. The source listing after the table explains the ROM algorithm; hardware may store those sites in ROM or dynamically generate equivalent sites after decoding the PTO-ISA header.\par\smallskip
}
\usepackage{enumitem}
\usepackage{pdflscape}
\setlist{nosep}
